% Síntesis axiomática incorporada desde el anexo de caminos.
% Las pruebas operatorias extensas permanecen en sus fuentes.

\section{Reconstruction des principes quantiques à partir de l’état généalogique}
\label{sec:postulados-cuanticos-genealogicos-rev2}
\index{mécanique quantique!reconstruction généalogique}
\index{état!réalisation de Hilbert}

La formulation hilbertienne organise la mécanique quantique au moyen d’états,
d’observables autoadjointes, d’une évolution unitaire, d’une composition tensorielle et d’une
règle probabiliste pour les résultats~\cite{vonNeumann1932}. HMT--MD
construit une couche antérieure à cette représentation : une section compatible,
sa route observable, la mémoire du transport et le lecteur qui publie un
résultat. L’espace de Hilbert n’est pas remplacé ; il apparaît comme codomaine d’une
réalisation linéaire déclarée.

La suite typée est
\[
 \boxed{\begin{aligned}
 \text{section généalogique}
 &\longrightarrow \text{réalisation linéaire}
 \longrightarrow \text{phase et superposition}\\
 &\longrightarrow \text{observable et évolution}
 \longrightarrow \text{probabilité et mesure}
 \end{aligned}}
\]
Chaque flèche conserve son domaine. En particulier, le vecteur d’état ne
remplace pas la section TPK, de même que la valeur réelle ne remplace pas
la généalogie numérique qui la produit.

\subsection{État généalogique et réalisation de Hilbert}

\begin{definition}[État généalogique réalisé]
Un état généalogique est un quintuplet
\begin{equation}
 \mathfrak G=(x_\infty,\Gamma,\mathcal F,\mathcal M,\mathcal R),
 \label{eq:estado-genealogico-realizado-rev2}
\end{equation}
où \(x_\infty\in\Omega_\infty\) est une section compatible,
\(\Gamma\) sa route observable, \(\mathcal F\) la famille de réalisations,
\(\mathcal M\) la mémoire des relèvements et \(\mathcal R\) le lecteur
déclaré. Soit \(\mathsf P_{\rm TPK}\) la catégorie dont les objets sont les
sections TPK préparées, finies ou coinductives, et dont les morphismes sont les
routes enregistrées compatibles avec leurs données de frontière. En particulier, la
restriction de profondeur est un morphisme
\[
 r_{n+1,n}:x_{n+1}\longrightarrow x_n.
\]
Soit \(\mathsf{Hilb}_{\mathbb R,J}\) la catégorie dont les objets
sont des paires \((\mathcal H_{\mathbb R},J)\), avec \(\mathcal H_{\mathbb R}\)
un espace de Hilbert réel et \(J\) orthogonal, \(J^2=-I\), et dont les morphismes sont
des applications linéaires réelles bornées \(T\) qui entrelacent les structures
complexes, \(TJ_1=J_2T\). Une réalisation quantique est un foncteur
\[
 \mathcal Q:\mathsf P_{\rm TPK}\longrightarrow
 \mathsf{Hilb}_{\mathbb R,J},
\]
et une \emph{réalisation pointée} est la paire \((\mathcal Q,\Psi)\) qui associe
à chaque section préparée \(x\) un vecteur non nul
\(\Psi_x\in\mathcal Q(x)\), compatible avec les restrictions du système
projectif.
\end{definition}

Avec la convention de linéarité dans le second argument, la forme hermitienne associée
s’écrit
\[
 \langle u,v\rangle_{\mathbb C}
 =\langle u,v\rangle_{\mathbb R}
  -i\langle u,J_E v\rangle_{\mathbb R}.
\]
Choisir la carte \(J_E\leftrightarrow i\) réalise la structure complexe qui, dans les Parties I--III, est construite comme racine interne de \(-I\). Pour des préfixes de profondeur \(n\), posons \(\mathcal H_n:=\mathcal Q(x_n)\) et
\[
 p_{n+1,n}:=\mathcal Q(r_{n+1,n}):
 \mathcal H_{n+1}\longrightarrow\mathcal H_n.
\]
La compatibilité entre restriction et réalisation exige
\begin{equation}
 p_{n+1,n}\Psi_{n+1}(x_{n+1})=\Psi_n(x_n).
 \label{eq:compatibilidad-realizacion-hilbert-rev2}
\end{equation}
Si \(r_{\infty,n}:x_\infty\to x_n\) est la restriction coinductive et
\(p_n:=\mathcal Q(r_{\infty,n}):\mathcal H_\infty\to\mathcal H_n\), une
limite \(\mathcal H_\infty\) réalise la section complète lorsque
\(p_n\Psi_\infty(x_\infty)=\Psi_n(x_n)\) à chaque profondeur.

\begin{principle}[État quantique réalisé]
L’état pur associé à une section est représenté par le rayon du vecteur non nul \(\Psi=\Psi_\infty(x_\infty)\) ; une préparation statistique est représentée par un opérateur densité positif de trace un sur la même réalisation. L’identification de la phase globale agit dans la représentation : la section et son registre demeurent dans le domaine généalogique.
\end{principle}

\subsection{Superposition et phase interne}

Si \(\Psi,\Phi\in\mathcal H_\infty\) et \(a=\alpha+i\beta\in\mathbb C\), l’action de \(a\) dans la présentation réelle est \(\alpha I+\beta J_E\). La combinaison \(a\Psi+b\Phi\) se forme après transport des deux possibilités vers une fibre commune. Les relèvements généalogiques de \cref{chap:genealogia-operaciones-rev2} prennent leurs valeurs dans \(\text{sortie}\times\text{mémoire}\) ; ils appartiennent au domaine TPK et ne se multiplient pas comme des opérateurs. Si une réalisation projective associe des opérateurs unitaires \(U_g\) aux transformations visibles, son défaut de composition est décrit séparément par
\[
 U_gU_h=c(g,h)U_{gh},
 \qquad c(g,h)\in U(1),
\]
avec \(c(e,g)=c(g,e)=1\) et
\[
 c(g,h)c(gh,k)=c(h,k)c(g,hk).
\]
Un cocycle à valeurs dans un autre groupe de mémoire n’agit dans la réalisation qu’après déclaration d’une représentation de ce groupe qui entrelace les opérateurs \(U_g\).

\begin{principle}[Superposition généalogique]
Les possibilités compatibles sont d’abord transportées dans \(\mathsf P_{\rm TPK}\), puis combinées linéairement dans \(\mathcal H_\infty\). L’action de \(J_E\) représente la phase relative et le poids du parcours conserve la donnée de relèvement déclarée.
\end{principle}

L’interférence s’exprime par les termes croisés
\[
 \|\Psi_1+\Psi_2\|^2
 =\|\Psi_1\|^2+\|\Psi_2\|^2
  +2\operatorname{Re}\langle\Psi_1,\Psi_2\rangle.
\]
Deux parcours ayant les mêmes extrémités peuvent donc contribuer avec des phases différentes sans que leur publication conjointe efface la différence des histoires.

\subsection{Observables et calcul spectral}

Soit \(O:\Omega_\infty\to Y\subseteq\mathbb R\) un lecteur généalogique réel. Une réalisation quantique de \(O\) est un opérateur autoadjoint \(\widehat O\), de domaine dense \(D(\widehat O)\), dont la PVM satisfait la condition de compatibilité
\[
 P_O(B)\Psi_x
 \quad\hbox{représente la restriction de \(x\) à \(O^{-1}(B)\)}
\]
pour tout borélien \(B\subseteq\mathbb R\) dans le support de la préparation. En dimension finie,
\[
 \widehat O=\sum_a o_a P_a,
 \qquad
 P_a P_b=\delta_{ab}P_a,
 \qquad
 \sum_a P_a=I;
\]
en dimension générale, la somme est remplacée par l’intégrale spectrale \(\widehat O=\int\lambda\,\mathrm dP_O(\lambda)\).

\begin{principle}[Observable typée]
L’observable physique est le couple formé par le lecteur dans le domaine généalogique et sa représentation autoadjointe. Le calcul spectral exprime les résultats possibles de la représentation sans effacer l’opération qui les a définis.
\end{principle}

Pour un état pur normalisé \(\Psi\in D(\widehat O^2)\),
\[
 \langle O\rangle_\Psi
 =\langle\Psi,\widehat O\Psi\rangle,
 \qquad
 (\Delta_\Psi O)^2
 =\langle\Psi,(\widehat O-\langle O\rangle_\Psi)^2\Psi\rangle.
\]
La formulation au moyen de \(\operatorname{tr}(\rho\widehat O)\) étend la définition aux états mixtes lorsque \(\rho\widehat O\) est à trace ; le deuxième moment exige également que \(\rho\widehat O^2\) soit à trace.

\subsection{Évolution, Hamiltonien et équation de Schrödinger}

L’évolution généalogique préserve la composition des transports. Sa
réalisation satisfait
\[
 U(t_3,t_2)U(t_2,t_1)=U(t_3,t_1),
 \qquad U(t,t)=I.
\]
Lorsqu’il préserve le produit scalaire et satisfait
\(U(t)J_E=J_EU(t)\), il est orthogonal dans la carte réelle et unitaire dans la carte
complexe. Pour un groupe unitaire à un paramètre fortement continu, le
théorème de Stone fournit un générateur autoadjoint et \(J_E\)-linéaire
\(H\)~\cite{Stone1932} :
\[
 U(t)=\exp\!\left(-J_E H t/\hbar_{\rm int}\right).
\]
Sur le domaine \(D(H)\), la dérivation temporelle donne
\begin{equation}
 \hbar_{\rm int}J_E\dot\Psi(t)=H\Psi(t),
 \label{eq:schrodinger-genealogico-rev2}
\end{equation}
et la carte \(J_E\leftrightarrow i\) produit
\(i\hbar_{\rm int}\dot\Psi=H\Psi\).

Le calendrier \(C_{108}\) réalise cette correspondance de manière finie :
le Hamiltonien autoadjoint construit dans sa source a pour
exponentielle, pendant \(t_0\), le décalage unitaire d’une itération. Le
groupe continu requiert en outre la continuité forte ; il ne se déduit pas
seulement de la périodicité du calendrier.

\subsection{Probabilité et représentation quadratique des branches}

Soit \(\mathcal C_n=\{C_s:s\in\mathcal S_n\}\) la partition de \(\Omega_\infty\) par un ensemble fini \(\mathcal S_n\) de préfixes de profondeur \(n\), et soit \(\mu\) une mesure cylindrique de probabilité compatible, \(\mu(\Omega_\infty)=1\). Soit \(E\) un espace de Hilbert réel muni d’une structure complexe \(J_E\). Dans \(\mathcal H_n=\ell^2(\mathcal S_n)\otimes E\), on définit
\begin{equation}
 \Psi_n=\sum_{s\in\mathcal S_n}
 \sqrt{\mu(C_s)}\,|s\rangle\otimes e^{\theta_s J_E}u_0,
 \label{eq:estado-amplitudes-medida-rev2}
\end{equation}
avec \(u_0\) unitaire.

\begin{theorem}[Représentation quadratique d’une mesure de branches]
\label{thm:born-cilindrico-rev2}
Pour un événement cylindrique \(A=\bigcup_{s\in I_A}C_s\), soit
\[
 P_A=\left(\sum_{s\in I_A}|s\rangle\langle s|\right)\otimes I_E.
\]
Alors
\[
 \|P_A\Psi_n\|^2=\mu(A).
\]
\end{theorem}

\begin{proof}
Les vecteurs de préfixe sont orthogonaux et \(e^{\theta_s J_E}\) préserve la norme. Par conséquent,
\[
 \|P_A\Psi_n\|^2
 =\sum_{s\in I_A}\mu(C_s)=\mu(A).
\]
\end{proof}

La règle quadratique est obtenue ici pour la PVM des événements cylindriques de profondeur finie comme représentation exacte de la mesure de branches déclarée~\cite{Born1926}. Réciproquement, le théorème de Gleason caractérise, en dimension au moins trois, les mesures additives non contextuelles sur les projecteurs au moyen d’opérateurs densité~\cite{Gleason1957}. Le couplage avec l’appareil et le conditionnement du résultat sont développés ensuite dans la théorie de la mesure.

\subsection{Incompatibilité des lecteurs et incertitude}

Les translations et phases nonadiques définissent le couple de Weyl
\[
 \mathcal H_9=\ell^2(\mathbb Z/9\mathbb Z),
 \qquad
 \omega_9=e^{2\pi J_E/9},
 \qquad
 X|n\rangle=|n+1\rangle,
 \qquad
 Z|n\rangle=\omega_9^n|n\rangle.
\]
Les indices se lisent modulo \(9\), et \(\omega_9\) est le scalaire complexe réalisé par la structure \(J_E\) sur \(\mathcal H_9\). L’identité
\[
 ZX=\omega_9XZ
\]
réalise le groupe fini de Heisenberg d’ordre \(9^3=729\). Ce cardinal appartient au groupe des phases et des translations et conserve son type par rapport aux autres occurrences de \(729\).

Pour des observables autoadjointes \(A\) et \(B\), un vecteur normalisé \(\Psi\in D(AB)\cap D(BA)\) ayant des deuxièmes moments finis permet d’appliquer Cauchy--Schwarz et d’obtenir l’inégalité de Robertson--Heisenberg
\begin{equation}
 \Delta_\Psi A\,\Delta_\Psi B
 \geq\frac12
 \bigl|\langle\Psi,[A,B]\Psi\rangle\bigr|.
 \label{eq:robertson-hmt-rev2}
\end{equation}
L’incompatibilité se présente ainsi à des niveaux coordonnés : le commutateur des évaluations d’APP, la relation de Weyl dans la représentation et la borne de dispersion pour les lectures physiques~\cite{Heisenberg1927,Robertson1929}.

Pour \(\Psi\neq0\), la transformée de Fourier discrète fournit la borne de support
\[
 |\operatorname{supp}\Psi|\,
 |\operatorname{supp}\widehat\Psi|\geq9,
\]
et, pour des bases mutuellement non biaisées de position et de phase, l’inégalité entropique de Shannon, avec des logarithmes naturels,
\[
 H_Q(\Psi)+H_P(\Psi)\geq\log9
\]
offre une troisième formulation exacte~\cite{MaassenUffink1988}. Ici, \(H_Q\) et \(H_P\) sont les entropies de Shannon des distributions de probabilité dans les deux bases.

\subsection{Produit tensoriel de parcours et génération de l’intrication par mémoire partagée}

La composition généalogique recolle des sections sur des données de frontière. Fixons sur \(\mathsf P_{\rm TPK}\) une structure monoïdale \((\boxtimes,\mathbf 1,a,l,r)\), avec produit, unité et cohérences d’associativité et d’unité déclarés. Une réalisation monoïdale forte fournit des isomorphismes naturels
\[
 \mu_{A,B}:
 \mathcal Q(A)\widehat\otimes\mathcal Q(B)
 \xrightarrow{\sim}\mathcal Q(A\boxtimes B).
\]
Dans l’image hilbertienne apparaît \(\mathcal H_A\otimes\mathcal H_B\). Un état pur non séparable admet une décomposition de Schmidt
\[
 |\Psi\rangle
 =\sum_r\sqrt{p_r}\,|r\rangle_A\otimes|r\rangle_B,
\]
et satisfait \(p_r\geq0\), \(\sum_rp_r=1\), avec \(\{|r\rangle_A\}\) et \(\{|r\rangle_B\}\) des bases orthonormales de Schmidt, et ses réductions possèdent le même spectre non nul. L’entropie associée à ce spectre reçoit, dans le chapitre de Barbero--Immirzi, une réalisation d’aire et de bord.

\begin{principle}[Composition généalogique]
Le produit tensoriel représente un recollement de sections sur une frontière partagée. L’intrication exprime la non-factorisation de la section réalisée relativement à la bipartition choisie.
\end{principle}

\subsection{Théorème de reconstruction cinématique}

\begin{theorem}[Reconstruction hilbertienne conditionnée de l’état généalogique]
\label{thm:cinematica-cuantica-tpk-rev2}
Soit \(\Omega_\infty\) un espace de sections compatibles, soit \((\mathsf P_{\rm TPK},\boxtimes,\mathbf 1)\) la catégorie monoïdale des parcours munie des cohérences déclarées, et soit
\[
 \mathcal Q:\mathsf P_{\rm TPK}\longrightarrow
 \mathsf{Hilb}_{\mathbb R,J}
\]
une réalisation monoïdale forte, pointée par une section compatible \(\Psi_x\in\mathcal Q(x)\). Supposons que :
\begin{enumerate}
 \item \(J_E^2=-I\), \(J_E\) préserve le produit scalaire et les morphismes de \(\mathcal Q\) entrelacent les structures complexes ;
 \item chaque lecteur réel déclaré est représenté par un opérateur autoadjoint dont la PVM satisfait la condition de compatibilité avec les fibres du lecteur ;
 \item chaque évolution continue est représentée par un groupe unitaire à un paramètre fortement continu \(U(t)\), avec \(U(t)J_E=J_EU(t)\), et un générateur autoadjoint \(J_E\)-linéaire \(H\) ;
 \item la préparation induit une mesure cylindrique de probabilité compatible et l’état \(\Psi_n\) de \eqref{eq:estado-amplitudes-medida-rev2}.
\end{enumerate}
Alors la carte complexe de \(\mathcal Q\) réalise des états sous forme de rayons ou de densités, la superposition complexe, des observables autoadjointes, une évolution unitaire engendrée par \(H\), la composition tensorielle et des probabilités quadratiques pour la PVM cylindrique de profondeur finie.
\end{theorem}

\begin{proof}
La première hypothèse transforme chaque fibre réelle en un espace de Hilbert complexe et rend les transports complexes linéaires. La deuxième permet d’appliquer le théorème spectral avec le lecteur typé. La troisième fournit le Hamiltonien par le théorème de Stone et l’équation différentielle sur \(D(H)\). La quatrième et \cref{thm:born-cilindrico-rev2} donnent la règle quadratique sur les événements cylindriques. La monoïdalité forte réalise le produit tensoriel et sa notion de séparabilité.
\end{proof}

\begin{statusbox}[title={Résultat et portée}]
La reconstruction hilbertienne est coordonnée comme propriété d’une réalisation de l’état généalogique : la structure complexe provient de \(J_E^2=-I\) ; l’évolution, de la représentation des parcours ; la probabilité, d’une mesure compatible ; et l’intrication, de la composition sur une frontière partagée. Le théorème est exact sous ses quatre hypothèses et n’identifie pas automatiquement TPK à une théorie quantique concrète. La construction d’une réalisation physique concrète conserve ses sources : Hamiltonien, appareil, lecteurs, domaine et conditions aux limites.
\end{statusbox}
